\begin{pastexam}[label={px:ntu-111-bime-ds-05}]{臺大生機所 111 學年度（C）第 5 題}
（問答）

(15\%) Please explain how a 2-3 tree achieves the proposition: every path from root to null link
has the same length.
\end{pastexam}

\begin{pastsoln}
null link 是葉節點底下空的子節點指標，所以「每條從根到 null link 的路徑一樣長」等價於「所有葉節點
都在同一層」，這正是 2-3 樹定義的第 3 條（資結書第~11.1~節）。 要說明的是：插入和刪除之後，這個性質
為什麼一直保持。 關鍵在 2-3 樹\textbf{從來不在底下長出新的一層，只在根長高、只在根變矮}。 對操作的
次數歸納，假設操作前所有葉節點都在同一層：

\textbf{插入}（資結書第~11.1.3~節）。 新鍵一定放進一個\textbf{已經存在}的葉節點，不會像二元搜尋樹那樣
在葉子底下掛一個新節點，所以沒有路徑變長。
\begin{itemize}
\item 葉節點原本有 $1$ 個鍵：變成 $2$ 個鍵，結構不變。
\item 原本有 $2$ 個鍵：變成 $3$ 個鍵，溢位，\textbf{分裂}成兩個各 $1$ 個鍵的節點，中間的鍵送進父節點。
兩個新節點都在原來那一層，父節點多一個子節點，深度都沒有改變。 父節點溢位就同樣往上分裂。
\item 分裂一路到根：根分裂成兩個節點，中間的鍵成為\textbf{新的根}放在最上面。 這時\textbf{每一條}路徑
都在最上面多了一步，長度同時加 $1$，仍然相等。
\end{itemize}

\textbf{刪除}（資結書第~11.1.4~節）。 內部節點的鍵先用後繼取代，所以真正拿掉鍵的地方一定是葉節點。
\begin{itemize}
\item 葉節點還剩 $1$ 個鍵：結構不變。
\item 葉節點變空：相鄰兄弟有 $2$ 個鍵就\textbf{借}一個（經過父節點轉一手），節點個數和位置都不變；兄弟只有
$1$ 個鍵就\textbf{合併}：空節點、父節點的分隔鍵、兄弟合成一個節點，放在\textbf{同一層}，父節點少一個子
節點。 父節點變空就同樣往上處理（內部節點借的時候，兄弟的一棵子樹跟著搬過來，那棵子樹的高度和
原本的子樹相同，所以葉子仍在同一層）。
\item 一路到根：根變空時把它拿掉，唯一的子節點成為新的根，\textbf{每一條}路徑同時少一步。
\end{itemize}

所以每次操作後，所有從根到 null link 的路徑長度要嘛都不變，要嘛同時加 $1$ 或同時減 $1$，永遠相等。
空樹和只有一個節點的樹一開始就成立，由歸納，任何一連串插入、刪除之後都成立。 這也是 2-3 樹高度只有
$O(\log n)$ 的原因：每個內部節點至少 $2$ 個子節點，所有葉子又在同一層，$n$ 個鍵的高度 $\le \lg(n+1)$。

對照：二元搜尋樹插入時在葉子底下新增節點，只有那一條路徑變長，所以會失去平衡；2-3 樹把「長高」
集中到根，一次讓所有路徑一起變長。
\end{pastsoln}
